\documentclass[11pt]{article}
\usepackage[T1]{fontenc}
\usepackage{lmodern,microtype,amsmath,amssymb,amsthm,mathtools}
\usepackage[a4paper,margin=24mm]{geometry}
\usepackage[dvipsnames]{xcolor}
\usepackage{enumitem,booktabs,longtable,array}
\usepackage[colorlinks=true,linkcolor=MidnightBlue,urlcolor=MidnightBlue]{hyperref}
\hypersetup{pdftitle={Local laws, delocalization, and weak rigidity: formalization guide},pdfsubject={Checked source theorems and explicit project validation status}}
\usepackage{fancyhdr}
\usepackage{listings}
\lstset{basicstyle=\ttfamily\footnotesize,breaklines=true,columns=fullflexible,keepspaces=true}
\setlist{nosep,leftmargin=1.5em}
\setlength{\parindent}{0pt}
\setlength{\parskip}{6pt}
\pagestyle{fancy}
\fancyhf{}
\fancyhead[L]{\small Formalization guide}
\fancyhead[R]{\small Source guide: 5 November 2026}
\fancyfoot[C]{\thepage}
\newcommand{\code}[1]{\texttt{\small\detokenize{#1}}}
\newcommand{\PP}{\mathbb P}
\newcommand{\EE}{\mathbb E}
\newcommand{\F}{\mathcal F}
\newcommand{\Om}{\Omega}
\newcommand{\im}{\operatorname{Im}}
\newcommand{\supp}{\operatorname{supp}}
% Flip only after the final integrated Lean project check succeeds.
\newif\ifIntegratedAudit
\IntegratedAudittrue
\newcommand{\auditstatus}{\ifIntegratedAudit Paper-form theorems and independent kernel replay checked\else Paper-form theorems checked; fresh kernel replay pending\fi}
\newcommand{\publicationstatus}{\ifIntegratedAudit
\textbf{Validation.} The independent integrated rebuild and trust-level-zero kernel replay have passed.
\else
\textbf{Publication gate.} The paper-form source declarations have compiled. The updated
integrated rebuild and trust-level-zero kernel replay must pass before this guide is published.
\fi}
\begin{document}
\begin{center}
  {\LARGE\bfseries Local laws, delocalization, and weak rigidity}\\[5pt]
  {\large A mathematical and Lean formalization guide}\\[5pt]
  {\small \auditstatus}
\end{center}

\colorbox{gray!10}{\parbox{\dimexpr\linewidth-2\fboxsep}{\small
\textbf{Checked paper-form statements.} Weak rigidity and all-vector delocalization use the
literal structural event and its actual conditional uniform measure. A necessary correction
to reference Definition 2.6 is stated below and proved consistent with the finite reference
class. The main spectral declarations have no analytic certificate hypotheses. \publicationstatus}}

\section{Objects, conditioning, and the scope of the reference class}

Fix an integer $d\geq3$. The probability space is the genuine uniform finite set of simple
$d$-regular graphs on $\{0,\ldots,N-1\}$. Ensemble nonemptiness is explicit in the source
theorems. Write $H=A/\sqrt{d-1}$ and $G(z)=(H-z)^{-1}$ for $\im z>0$.
The paper's real-radius structural event is
\[
 \begin{gathered}
 \F_N(c,\omega)=\left\{\max_v\operatorname{excess}(B_R(v))\leq\omega,
 \quad\#\{v:B_R(v)\text{ is not a tree}\}\leq N^c\right\},\\[2pt]
 R=\frac{c\log N}{4\log(d-1)},\qquad 0<c<1.
 \end{gathered}
\]
Because graph distances and vertex counts are integers, the real-radius balls and the
real count bound are proved exactly equivalent to the integer-floor definitions used by
the local-law development. The final probability statement is under the actual uniform
law on the structural subtype: for every fixed real $C>0$, eventually in admissible $N$,
\[
 \PP_{\operatorname{Unif}(\F_N)}(\text{failure})
 =\frac{\PP(\F_N\cap\{\text{failure}\})}{\PP(\F_N)}\leq N^{-C}.
 \tag{1}
\]
The event has mass at least $1/2$ eventually, so its nonemptiness and the denominator's
positivity are proved. The final wrappers require no nonemptiness assumption on this
structural subtype beyond eventual nonemptiness of the original regular-graph ensemble.

\textbf{Necessary correction to Definition 2.6.}
Its printed estimate (2.10), $|G_{ii}(z)|\gtrsim1$ for \emph{all} $z\in\mathbb C_+$, is impossible:
the genuine self-adjoint resolvent satisfies $|G_{ii}(\mathrm iY)|\leq1/Y$.
The theorem \code{no_global_positive_diagonal_lower_bound} proves this obstruction.
The global reference class therefore uses
$|m_d(z)|/C_d\leq |G_{ii}(z)|$; it recovers constant lower and upper diagonal bounds on
every fixed bounded spectral region. This correction is explicit, rather than a claim that
the inconsistent printed definition was formalized unchanged.

The corrected \code{PaperReferenceThreshold} quantifies all connected, possibly infinite
cores with finite deficit support and an actual finite feedback-edge set whose deletion
leaves a forest. Their deficit sum plus feedback-edge count is at most $\omega$.
Its Green function is the genuine $\ell^2$ resolvent of the literal infinite tree extension.
Finite induced approximations, neighbor-closed embeddings, and strong resolvent convergence
prove that this all-core class is equivalent to \code{FiniteReferenceThreshold}.
The positive integer \code{paperOmega d} is independently defined as the maximum of this
all-core class; the equality \code{paperOmega_eq_finiteOmega} is then proved.
The final structural event uses this maximum, with no reference-class hypothesis left over.

\begin{lstlisting}
SpectralConsequences.paperReferenceThreshold_iff_finite
SpectralConsequences.paperOmega_eq_finiteOmega
SpectralConsequences.allCore_reference_bounded_diagonal
SpectralConsequences.paperStructural_iff
SpectralConsequences.conditional_uniform_eq_subtype_uniform
\end{lstlisting}
The semicircle and Kesten--McKay density Stieltjes integrals are also identified with the
actual scalar roots, rather than postulated.

\section{Theorem 4.2 supplies the source local law}

The existing \code{LocalLaw} project is the foundation. Its main declaration is
\begin{lstlisting}
LocalLaw.source_theorem_4_2
LocalLaw.source_theorem_4_2_of_reference_threshold
\end{lstlisting}
It permits real $a,b$ with $a\geq12$ and $b\geq25a$, and $0<\beta\leq c/32$.
The conclusion is uniform over the upper-half-plane domain
$\im z\geq (\log N)^b/N$, with the actual nominal model radius
$r_{\rm nom}=\lfloor\beta\log N/\log(d-1)\rfloor$ and the article's actual error.
All first- and second-stage graph resampling estimates used by this theorem are proved.
No local-law, concentration, model-comparison, or resampling certificate is a hypothesis of
its primary source declaration.

The proof uses a shorter internal model radius, losing $O(\log\log N)$ layers.
It keeps the nominal error fixed, avoiding comparisons across the discontinuous
error-inflation branches.

\section{Theorem 1.4: all eigenvectors from Ward's identity}

For a Hermitian matrix and an arbitrary eigenpair $Hv=\lambda v$, the proved deterministic
criterion is
\[
 |v_i|^2\leq \eta\,\im G_{ii}(\lambda+\mathrm i\eta)
                  \sum_j|v_j|^2,\qquad \eta>0.
 \tag{2}
\]
It includes the zero vector and does not assume a chosen eigenbasis or a simple eigenvalue.
The source local law gives a degree-dependent uniform diagonal cap. With $a=12$, $b=300$,
and $\eta=(\log N)^{300}/N$, this yields simultaneously for every eigenpair
\[
 |v_i|^2\leq C_d\frac{(\log N)^{300}}{N}\sum_j|v_j|^2 .
 \tag{3}
\]
Consequently, every normalized eigenvector satisfies
$\|v\|_\infty\leq\sqrt{C_d}(\log N)^{150}/\sqrt N$.
Absorbing the fixed degree-dependent constant into one more logarithm gives the literal
homogeneous conclusion
\[
 \|v\|_\infty\leq\frac{(\log N)^{151}}{\sqrt N}\,\|v\|_2 .
\]
The final predicate quantifies \emph{all complex eigenvalues and all complex vectors};
Hermitian reality is proved for nonzero eigenvectors, and the zero vector is handled
separately. No normalization, chosen eigenbasis, or simple-eigenvalue hypothesis is needed.
The conclusion satisfies (1) on the literal structural subtype at \code{paperOmega d}.
Its canonical wrapper takes only $d$, $c$, their numerical inequalities, and eventual
ensemble nonemptiness; the auxiliary $\beta$ is selected internally.

\begin{lstlisting}
SpectralConsequences.eigenvector_coordinate_sq_le_green_im_mul_sum
SpectralConsequences.paperAllEigenvectors_iff_paperDelocalized
SpectralConsequences.full_paper_theorem_1_4
SpectralConsequences.full_paper_theorem_1_4_ratio
\end{lstlisting}
The deterministic criterion is in \code{EigenvectorWard.lean}. The source local-law
application, constant absorption, unrestricted complex-eigenpair statement, and literal
conditional uniform law are connected in the subsequent paper-form wrappers.

\section{Counting, quantiles, and the distinction from confinement}

The sorted eigenvalues are in decreasing order, with multiplicities retained.
The constant vector supplies the distinguished index $0$ eigenvalue
$\tau_d=d/\sqrt{d-1}$. The other indices remain in the empirical measure even if their
value also equals $\tau_d$.

The actual Kesten--McKay density is normalized to mass one. Its compact CDF is continuous,
monotone, and strictly increasing in the interior. The reference-mass estimate
\[
 F_d(y)-F_d(x)\geq c_*|y-x|^{3/2},\qquad -2\leq x\leq y\leq2,
\]
provides an inverse $2/3$ modulus for its quantiles. Poisson smoothing converts the trace
part of Theorem 4.2 into a uniform empirical-CDF bound; one checked source consequence is
\[
 \sup_E|F_N(E)-F_d(E)|\leq N^{-\beta/8}
\]
with every polynomial conditional failure power.

A CDF estimate alone controls clamped eigenvalues, not outliers. The deterministic rigidity
step explicitly also requires confinement. In a general form, CDF error $\delta$ and
confinement error $e$ imply
\[
 |\lambda_i-\gamma_i|\leq e+(\delta/c_*)^{2/3}.
 \tag{4}
\]
In zero-based Lean indexing, the classical locations satisfy
$F_d(\gamma_i)=(N-1-i)/N$. These are the exact source index levels after translating
the article's one-based indexing. For one-based $2\leq i\leq N$, the exact identity
$\int_{\gamma_i}^{2}\rho_d(x)\,dx=i/N$ is proved, including $\gamma_N=-2$.
The source confinement event excludes only index $0$. It does not discard every occurrence
of the regular eigenvalue or assume connectedness.

\begin{lstlisting}
SpectralConsequences.source_counting_rate
SpectralConsequences.empiricalCDF_error_of_trace_sqrt
SpectralConsequences.SquareRootCDF.inverse_modulus
SpectralConsequences.SquareRootCDF.rigidity_of_cdf_and_confinement
SpectralConsequences.RegularGraph.correctedSourceTrace_eq_positive_sum
\end{lstlisting}
The generic inverse-modulus and confinement-to-rigidity lemmas are in
\code{QuantileRigidity.lean}; the source counting wrapper is in
\code{SourceCountingProbability.lean}.

\section{Section 7: physical fluctuations and reversible moments}

\subsection{The observables and the trivial mode}

Let $q=Q(G,z)$ be the physical averaged oriented cavity observable, let $m=m_{\rm sc}(z)$,
and let $m_d$ be the Kesten--McKay transform. The depth $\ell$ is selected by the actual
nonresonance rule. The finite cavity and full-root references are
$Y_\ell(q)=\operatorname{cavityIterate}(z,q,\ell+1)$ and $X_\ell(q)$.
The two corrected physical residuals are
\[
 F_Q=q-Y_\ell(q)-\delta_Q,\qquad
 F_m=m_N-X_\ell(q)-\delta_m .
\]
The corrections are explicit physical tree-weight multiples of
$\pi_N(z)=N^{-1}(\tau_d-z)^{-1}$. Tree root-to-boundary columns, the exact number of
boundary labels, and the active labels after deleting a root neighbor are all proved for
the actual graph balls. Thus $\delta_Q$ and $\delta_m$ are identified with genuine
boundary-weight sums, including their multiplicities and zero inactive branches.

Subtracting the one distinguished pole gives a positive spectral sum over all indices
other than $0$. This is needed to exclude extra trivial eigenvalues as well as ordinary
outliers. The scalar linearization and the two corrections obey a proved global cancellation
estimate, including near $\tau_d$; a pole-separation assumption is not imposed.

\begin{lstlisting}
SpectralConsequences.actual_depth_trivial_correction_cancellation
SpectralConsequences.eventually_source_pole_scale_quadratic
SpectralConsequences.RegularGraph.cavity_boundary_all_label_weights
SpectralConsequences.RegularGraph.full_boundary_offdiag_correction
\end{lstlisting}
In particular $\Psi(z)|\pi_N(z)|\leq C_{d,B}\varepsilon(z)^2$ on a bounded spectral domain.
Near $\tau_d$, the edge distance is positive; away from $\tau_d$, the fractional error
floor controls the inverse sample size. This turns the q-to-physical correction change
into a quadratic drift cost.

\subsection{Sampling geometry and exits}

The resampling probability space includes both the uniformly random boundary labeling
and the genuine independent edge-sampling fiber. Its switching map is reversible.
Large protected tree patches imply actual acceptance of every label and preserve the
same-radius excess bound and the root tree event. These are deterministic statements
about the literal resampled graph, not assumed good-gadget conclusions.

At the smaller work parameter, two separate graph-count auxiliaries are used. The larger
radius count bounds the failure of protected sampling; the smaller count preserves the
original work bad-center threshold after switching. Both auxiliary complements have every
polynomial graph-probability tail. The conditional geometric exit probability is at most
$N^{-9/32}$. The auxiliary exponents are not used as a replacement for the target structural
bad-count exponent.

\begin{lstlisting}
SpectralConsequences.RegularGraph.eventually_source_geometric_switch_exit_concrete
SpectralConsequences.LabelledResamplingState.ProtectedPatches.same_radius_excess
SpectralConsequences.LabelledResamplingState.ProtectedPatches.same_radius_root_tree
\end{lstlisting}

\subsection{What enters the genuine moment recurrence}

The high-moment argument uses the complex test
$T_p(F)=F^{p-1}\overline F^{\,p}$, so that $F T_p(F)=|F|^{2p}$.
Exact finite disintegration lifts actual conditional entry means to the labelled-ball
kernel. The resulting source drift integral for either observable is bounded by
\[
 C(\log N)^{6a+7}\varepsilon^2\,
      \EE[\mathbf1_{\Om}|F|^{2p-1}]
 +C M^{2p-1}N^{-D} .
 \tag{5}
\]
The second term is a proved auxiliary-event loss. Uniform switched local caps are derived
from the old event for \emph{every} sampling assignment, so exceptional samples do not
require an unproved new-graph resolvent cap.

On both old and new source events, genuine Ward estimates bound the corrected global
increments by
\[
 V=C_{d,B}(\log N)^{2a}
        \frac{\im m_d+\varepsilon+\Psi}{N\eta}.
 \tag{6}
\]
The same variance satisfies $0\leq V\leq\varepsilon/(\log N)^{16}\leq\varepsilon$.
Actual local observable increments are bounded by a fixed logarithmic multiple of
$\varepsilon$. The exit integral is controlled separately, using the polynomial geometric
exit and all-power weak-event/auxiliary failures. The recurrence keeps the mixed
$\varepsilon_{\rm loc}V$ scale; it does not replace it by a first-order local-law error.
This proved Section 7 generalization retains fixed logarithmic losses, including
$(\log N)^{6a+7}\varepsilon^2$. They are absorbed into $N^e$ in the weak-rigidity argument.
It is not a claim of the printed log-free sharp Proposition 7.5 estimate.

\begin{lstlisting}
SpectralConsequences.global_reversible_moment_recurrence
SpectralConsequences.RegularGraph.eventually_source_moment_drift_integral
SpectralConsequences.RegularGraph.eventually_source_ward_variance_bound
SpectralConsequences.RegularGraph.eventually_source_ward_variance_small
SpectralConsequences.source_actual_high_moment_bounds
\end{lstlisting}
The actual source high-moment family, its physical-tail bridge, and the final rigidity
wrapper are checked. The reversible adapter requires increment bounds only on the nonzero
support of the actual sampling law; zero-weight auxiliary states do not introduce a spurious
hypothesis. No intermediate moment or residual-tail predicate remains as a hypothesis of
the final main theorem. \ifIntegratedAudit The independent integrated rebuild and
trust-level-zero kernel replay have also passed.\else The fresh integrated rebuild and
trust-level-zero kernel replay remain the publication gate.\fi

\section{The final weak-rigidity source statement}

For the original conditioning exponent $c$, set
\[
 c_{\rm aux}=\min(c/2,1/4),\qquad
 \beta=\frac{c_{\rm aux}}{256},\qquad
 \alpha=\frac{\beta}{32}=\frac{\min(c/2,1/4)}{8192}>0.
\]
The source reduction uses $a=12$, $b=300$. Confinement from the corrected physical trace,
combined with the checked counting law and (4), gives the weak-rigidity conclusion
\[
 \max_{1\leq i<N}|\lambda_i-\gamma_i|\leq N^{-\alpha}
\]
with the all-polynomial conditional bound (1). This is a weak positive-power rate, not the
optimal edge $N^{-2/3}$ scale. The original conditioning event and its original radius
remain in the final denominator, although the proof works internally at smaller radii.

The final public declarations are
\begin{lstlisting}
SpectralConsequences.full_paper_theorem_1_2
SpectralConsequences.full_paper_theorem_1_2_ratio
SpectralConsequences.paper_classical_location_tail_integral
\end{lstlisting}
They are in \code{FullPaperTheorems.lean} and have only the numerical hypotheses
$d\geq3$, $0<c<1$, followed by eventual ensemble nonemptiness. They explicitly supply
$\alpha>0$, retain multiplicities, and use the literal one-based range $2\leq i\leq N$.
The primary form uses uniform measure on \code{PaperGraph}; the ratio companion is exactly
equal to it by the proved conditioning identity. Earlier declarations such as
\code{canonical_source_rigidity_of_residual_tails} remain useful intermediate composition
lemmas, but their residual-tail hypothesis is discharged in the final theorem.

\section{How to check and navigate the Lean development}
\begingroup\small\setlength{\parskip}{4pt}

The root directories \code{LocalLaw/} and \code{SpectralConsequences/} contain the
local-law foundation and the spectral developments.\\
Final statements: \code{FullPaperTheorems.lean}.\\
Corrected all-core maximum: \code{PaperReferenceThreshold.lean}.\\
Each library's \code{Main.lean} and \code{Audit.lean} register its sources and query
principal theorem axioms.\\
Ward: \code{EigenvectorWard.lean}; counting: \code{CDFComparison.lean};\\
Section 7: \code{SourceMomentObservables.lean}.

From the repository root, the full check is
\begin{lstlisting}
lake exe cache get
lake env python3 scripts/check.py --kernel
\end{lstlisting}
With compiled dependencies installed, \code{./check.sh --kernel} runs the same checker.\\
See \code{docs/BUILDING.md} for dependency overrides and offline use. Under \code{docs/},
\code{public-statements.txt}, \code{proof-map.txt}, and \code{paper-match-audit.txt}
give the exact types, proof map, and source comparison.

The pinned Lean toolchain is 4.19.0 with its compatible Mathlib cache. The claimed
source declarations use only \code{propext}, \code{Classical.choice}, and \code{Quot.sound}.
The checker replays safe project declarations into a Mathlib base environment at kernel
trust level zero. Lean's official replay excludes unsafe/partial compiler auxiliaries;
all theorem proof terms are included in the safe mathematical declarations checked.
There are no custom axioms, \code{sorry} placeholders, or assumed source mean/drift/exit
certificates in these completed components. The source files and their audited declarations,
rather than an informal proof outline, determine the exact checked statements.
\endgroup

\end{document}
